
\documentclass[aspectratio=169]{beamer}

\mode<presentation>
\usetheme{Boadilla}
\definecolor{NTHU1}{RGB}{127,16,132}
\definecolor{NTHU2}{RGB}{228,210,231}
\definecolor{blue}{RGB}{30,90,205}
\definecolor{red}{RGB}{213,94,0}
\definecolor{green}{RGB}{0,128,0}
\definecolor{charcoalgray}{RGB}{108,104,104}
\setbeamercolor{title}{fg=NTHU1}
\setbeamercolor{frametitle}{fg=NTHU1}
\setbeamercolor{block title}{bg=NTHU1, fg=white}
\setbeamercolor{block body}{bg=NTHU2}
\setbeamercolor{structure}{fg=NTHU1}
\setbeamercolor{item projected}{fg=white}
\setbeamercolor{item}{fg=NTHU1}
\setbeamercolor{subitem}{fg=NTHU1}
\setbeamercolor{section in toc}{fg=NTHU1}
\setbeamercolor{description item}{fg=NTHU1}
\setbeamercolor{caption name}{fg=NTHU1}
\setbeamercolor{button}{bg=NTHU1, fg=white}
\setbeamercolor{caption name}{fg=NTHU1}
\usepackage{graphics}
\usepackage{tikz}
\usepackage{amsmath}
\usepackage{bbm}
\usetikzlibrary{decorations.pathreplacing}
\usepackage{geometry}
\usepackage{booktabs}
\usepackage{bookmark}
\usepackage{multirow, makecell}
\usepackage{float}
\usepackage{fancyvrb}
\usepackage{caption}
\captionsetup[figure]{singlelinecheck = false, format= hang, justification=centering}
\captionsetup[subfigure]{singlelinecheck = false, format= hang, justification=raggedright}
\usepackage{subcaption}
\usepackage{adjustbox}
\usepackage{hyperref}
\usepackage{threeparttable}
\usepackage{appendixnumberbeamer} 
\usepackage[T1]{fontenc}
\usepackage[default]{lato} 
\usepackage{smartdiagram}
\smartdiagramset{
  back arrow disabled = true,
  module minimum width=1cm,
  module x sep = 3,
  uniform arrow color=true,
}
\usepackage{xcolor}
\usepackage{colortbl}
  \definecolor{light-gray}{gray}{0.99}

\newenvironment{wideitemize}{\itemize\addtolength{\itemsep}{10pt}}{\enditemize}
\newenvironment{wideenumerate}{\enumerate\addtolength{\itemsep}{10pt}}{\endenumerate}
\newenvironment{widedescription}{\description\addtolength{\itemsep}{10pt}}{\enddescription}
\hypersetup{
colorlinks=true,
linkcolor=NTHU1,
filecolor=green, 
urlcolor=blue,
}
\beamertemplatenavigationsymbolsempty
\setbeamercolor{author in head/foot}{bg=white, fg=NTHU1}
\setbeamercolor{title in head/foot}{bg=white, fg=NTHU1}
\setbeamercolor{date in head/foot}{bg=white, fg=NTHU1}
\setbeamercolor{section in head/foot}{bg=white, fg=NTHU1}
\setbeamercolor{headline}{bg=white}
\setbeamertemplate{footline}{
    \leavevmode%
    \hbox{%
        \begin{beamercolorbox}[wd=.333333\paperwidth,ht=2.25ex,dp=1ex,center]{date in head/foot}%
            \usebeamerfont{date in head/foot}%\insertdate
        \end{beamercolorbox}%
        \begin{beamercolorbox}[wd=.333333\paperwidth,ht=2.25ex,dp=1ex,center]{title in head/foot}%
            \usebeamerfont{title in head/foot}\insertshorttitle
        \end{beamercolorbox}%
        \begin{beamercolorbox}[wd=.333333\paperwidth,ht=2.25ex,dp=1ex,right]{date in head/foot}%
            \usebeamerfont{date in head/foot}\insertframenumber{}\hspace*{2em}
        \end{beamercolorbox}}%
        \vskip0pt%
    }
    %% May need to script the introduction!!!!!!
%\setbeamercolor{page number in head/foot}{fg=NTHU1}
\setbeamertemplate{section in toc}[sections numbered]
\setbeamertemplate{subsection in toc}{\leavevmode\leftskip=3em\rlap{\hskip-1.75em\inserttocsectionnumber.\inserttocsubsectionnumber}
\inserttocsubsection\par}
\setbeamerfont{subsection in toc}{size=\footnotesize}
%\setbeamertemplate{headline}{%
  %\begin{beamercolorbox}[ht=5.5ex]{section in head/foot}
    %\vskip2pt\insertnavigation{0.33\paperwidth}\vskip2pt
  %\end{beamercolorbox}%
%}
\newenvironment{transitionframe}{\setbeamercolor{background canvas}{bg=white}\setbeamertemplate{footline}{} \begin{frame}[noframenumbering]}{\end{frame}}

\makeatletter
\let\@@magyar@captionfix\relax
\makeatother

\title[]{Conflict and Peace} % Change this regularly
\subtitle[]{Week 11: Consequences of migration and host society perception}
\author[Lee]{Seung-hun Lee }
\institute[]{Taipei School of Economics\\ National Tsing Hua University}

\date[]{November 14th, 2025}

\begin{document}
\begin{transitionframe}
\titlepage
\end{transitionframe}

%%% Color slides for section headers: Use for colloquium version (The ones bewteen \iffals and \fi)

\section{Overview of the topic}
\begin{frame}
\frametitle{What we know about migration so far} 
Share of migrants increasing across many countries
\begin{itemize}\setlength\itemsep{3pt}
\item 2024: 304 million people globally are migrants, roughly 4\% (UNDESA)
\item Global South to North migration: Seeking labor and educational opportunities
\item Rising share of migration across countries within Global South
\item Internally displaced persons on the rise: 83 million people (UNDESA)
\begin{itemize}
\item Those forced to flee violence, disasters, or crises but remain within country borders
\end{itemize}
\end{itemize}\medskip
Reasons for migration are also diverse
\begin{itemize}\setlength\itemsep{3pt}
\item Primary motivation are for labor and educational opportunities
\item But there are also political migrants: Those fleeing conflicts and political instability
\item Increasing share of those fleeing environmental disasters
\end{itemize}\medskip
With rising conflicts and natural disasters globally, understanding forced migration is crucial
\end{frame}

\begin{frame}
\frametitle{Growth in migrant populations, even in Asia} 
\begin{figure}
  \includegraphics[keepaspectratio, width=0.6\textwidth]{fig1.png}
\end{figure}
\begin{scriptsize}
Source: UN Department of Economics and Social Affairs (Note: For US, 9.2\% $\to$ 15.2\%; Europe: 7.1\% $\to$ 12.6\%)
\end{scriptsize}
\end{frame}

\begin{frame}
\frametitle{Migration dynamics and rationale} 
\begin{figure}
  \begin{subfigure}{0.5\textwidth}
  \includegraphics[keepaspectratio, width=\textwidth]{fig2.png}  
  \end{subfigure}
  \begin{subfigure}{0.45\textwidth}
  \includegraphics[keepaspectratio, width=\textwidth]{fig3.png}  
  \end{subfigure}
\end{figure}
\end{frame}

\begin{frame}
\frametitle{How do migrants integrate into host societies?} 
This is a two-way effort: Effort of migrants and the natives
\begin{itemize}\setlength\itemsep{3pt}
\item So far, emphasis on the Effects of migrants on host society
\begin{itemize}\setlength\itemsep{3pt}
\item Impact of migrant workers on themselves and natives (Peri et al 2016 and many more)
\item Impact of migrant voters on electoral outcomes (Barbsai et al 2017; Mayda et al 2022)
\item Migrant \& Second generation (Bleakly, Chin 2010; Bratsberg et al 2020)
\end{itemize}
\item We know relatively little about how natives in the host societies affect integration
\begin{itemize}\setlength\itemsep{3pt}
\item Are they perceived with prejudice as outsiders? How?
\item Are these based on accurate or misguided information?
\item What does it imply for migrant integration and policies?
\end{itemize}
\end{itemize}\medskip
In this lecture, we focus the impacts of forced migration (political and climate migrants)\\ \medskip
In addition, we look at how migrant perceptions of natives affect integration and policies
\end{frame}

\begin{transitionframe}
\frametitle{Roadmap for today}
\tableofcontents
\end{transitionframe}

\section{Consequences of migration}


\subsection{Card (1990): The Impact of Mariel Boatlift on the Miami Labor Market}
\begin{transitionframe}
\begin{center}
  \begin{huge}
    \textcolor{NTHU1}{%
Card (1990): \\ The Impact of Mariel Boatlift on the Miami Labor Market
    }
  \end{huge}
\end{center}
\end{transitionframe}

\begin{frame}
\frametitle{Q: How does migrant inflow affect labor market?}
Migration policy debates center around how migrants affect native workers' outcomes
\begin{itemize}\setlength\itemsep{3pt}
\item How does it affect their wages and employment?
\item If migrants substitute native workers, lower wages and employment for natives
\item If migrants complement native workers, effects may not be straightforward
\begin{itemize}\setlength\itemsep{3pt}
\item Migrants and natives may have different skill sets: May not compete for same jobs
\item Some natives (e.g. high-skilled) are probably not as affected as others (e.g. low-skilled)
\item Host society can absorb these changes (Demand for specific workers $\Uparrow$)
\end{itemize}
\item Still debating on how this affects native worker outcomes in various settings
\end{itemize}\medskip
We look at one of the earliest works that started this debate
\begin{itemize}\setlength\itemsep{3pt}
\item Leverages sudden and large inflow of Cuban migrants to a concentrated area (Miami)
\item Micro-data that identifies migrants \& pre-migrant labor market situations
\end{itemize}
\end{frame}

\begin{frame}
\frametitle{Migrant inflow following Mariel Boatlift (Exodo del Mariel)} 
What is Mariel Boatlift?
\begin{itemize}\setlength\itemsep{3pt}
\item Approx. 120,000 Cubans entered US in 05/1980 - 09/1980
\item After Cuban govt announcement that anyone who wanted to leave Cuba could do so
\item Half of them settled in Miami \& high \% of less-skilled workers with low English fluency
\item Lower occupation attainment when compared to earlier cohorts of Cuban migrants
\end{itemize}\medskip
Labor market in Miami around Mariel Boatlift
\begin{itemize}\setlength\itemsep{3pt}
\item Miami: Migrant intensive even before the Boatlift (35.5\% of population foreign-born)
\item Also had large Black American contingent (17.3\% of population)
\item Cubans and other Hispanics primarily work as craftsmen and operatives
\item Black Americans in Laborer and service-related occupations 
\item White Americans in in Technical and Manager positions 
\end{itemize}\medskip
\end{frame}

\begin{frame}
\frametitle{Data and empirical strategy} 
Data sources
\begin{itemize}\setlength\itemsep{3pt}
\item 1979-1985 individual data from Current population survey (CPS) 
\begin{itemize}\setlength\itemsep{3pt}
  \item Also includes pre-Boatlift Miami labor market outcomes
  \item Cubans separately identified in CPS (vs `other Hispanics')
  \item Relatively large Miami sample - 1,200 per month
  \end{itemize}
\end{itemize}\medskip
Empirical findings leverage pre vs post Boatlift trends in different cities
\begin{itemize}\setlength\itemsep{3pt}
\item Compare: Atlanta, LA, Houston, Tampa (similar demographics and economic growth)
\item Looks at hourly earnings, and unemployment across different demographic groups
\item Similar to DiD (but not exact - Miami had lower earnings than others pre-Boatlift)
\item Control for education, age, gender, race, gender $\to$ Generate predicted wage
\end{itemize}\medskip
\end{frame}

\begin{frame}
\frametitle{Wage differences stable over the period for non-Cubans} 
\begin{figure}
\includegraphics[keepaspectratio, width=0.8\textwidth]{c1990_t1.png}
\end{figure}
\end{frame}

\begin{frame}
\frametitle{Stable employment differentials for low-skilled Blacks} 
\begin{figure}
\includegraphics[keepaspectratio, width=0.6\textwidth]{c1990_t2.png}
\end{figure}
\begin{itemize}
  \item Differences likely due to 1982 recession (not Boatlift)
\end{itemize}
\end{frame}

\begin{frame}
\frametitle{Labor market outcomes for Cuban population stable across the period} 
\begin{figure}
\includegraphics[keepaspectratio, width=0.67\textwidth]{c1990_t3.png}
\end{figure}
\end{frame}

\begin{frame}
\frametitle{Subsequent studies and remaining questions} 
What did the subsequent studies find?
\begin{itemize}\setlength\itemsep{3pt}
\item Card (1990): No effect on non-Cuban and Cubans (Absorbed additional labor force)
\item Borjas (many, including 2017): Large wage $\Downarrow$ for high school dropouts
\item Peri \& Yasenov (2019): Using synthetic controls, finds limited effect on that high-school dropout wages
\item Clemens \& Hunt (2019): Wage decline explained by race composition changes uniquely within Miami CPS subsamples. Using stable samples, they find no effect on low-skill wages 
\end{itemize}\medskip
Remaining questions
\begin{itemize}\setlength\itemsep{3pt}
\item What happens to home countries of migrants?
\begin{itemize}\setlength\itemsep{3pt}
  \item Home country loses human capital that could have contributed
  \item Especially since most migrants are in the productive stages of their lives
\end{itemize}
\item How do migrants determine where to settle?
\begin{itemize}\setlength\itemsep{3pt}
  \item Many migrants can choose where to live, do they follow network of previous migrants?
\end{itemize}
\end{itemize}\medskip
\end{frame}

\subsection{Waldinger (2010): Quality Matters: The Expulsion of Professors and Consequences for PhD Student Outcomes in Nazi Germany}
\begin{transitionframe}
\begin{center}
  \begin{huge}
    \textcolor{NTHU1}{%
Waldinger (2010): \\ Quality Matters: The Expulsion of Professors and Consequences for PhD Student Outcomes in Nazi Germany
    }
  \end{huge}
\end{center}
\end{transitionframe}

\begin{frame}
\frametitle{Q: How does migration affect innovation at home nation?} 
Quality of workers determines potential for innovation
\begin{itemize}\setlength\itemsep{3pt}
\item University quality $\to$ successful professional career, innovation of students
\begin{itemize}
\item For PhD students, this means getting publications, patents, and future jobs
\end{itemize}
\item Many key determinants: Professors, department resources etc
\item Causal analysis difficult
\begin{itemize}\setlength\itemsep{3pt}
\item Selection bias: Top quality students select into top schools
\item Omitted variable bias: quality of laboratories and research programs
\end{itemize}
\end{itemize}\medskip
This paper leverages dismissal of professors to study how  migration affects innovation
\begin{itemize}\setlength\itemsep{3pt}
\item Nazi Germany dismissed `politically unreliable' professors from German Universities
\item Mathematics were particularly hard-hit: Many were Jews or were labeled as such
\item Looks at how sudden dismissal affected student outcomes and innovation
\end{itemize}
\end{frame}

\begin{frame}
\frametitle{Expulsion of professors by Nazi Regime} 
Who were dismissed by the Nazi government
\begin{itemize}\setlength\itemsep{3pt}
\item `Law for the Restoration of the Professional Civil Service'
\item Non-Aryans (Jews in particular), Any professors who supported opposition
\item Led to dismissals and `early retirements' from German Universities
\item Mathematics: About 18\% of all professors dismissed in 1933-34
\item Many fled to foreign universities (mostly US) and unable to advise former students
\item Top departments (Gottingen, Berlin) were hard hit
\end{itemize}\medskip
Who were the PhD students of Math in Germany
\begin{itemize}\setlength\itemsep{3pt}
\item Focus on PhD students who obtained degree by 1938
\item Those from top universities are more likely to publish and become full professors
\item Some also got high school teacher license
\end{itemize}\medskip
\end{frame}

\begin{frame}
\frametitle{Data and empirical strategy} 
Data sources
\begin{itemize}\setlength\itemsep{3pt}
\item Students: Career outcomes of `23-`38 cohort from German Mathematical Association
\item Professor data: List of Displaced German Scholars, University Calendar
\item Department quality: Avg \# of citation-weighted publications in top journals
\end{itemize}\medskip
Regression that leverages exogenous changes in department quality due to dismissals
\[
Y_{idt}=\beta_1+\beta_2\text{Quality}_{dt-1}+\beta_3\text{S/F Ratio}_{dt-1}+\gamma X_{idt}+ \phi_t + \psi_d + e_{idt}
\]
\vspace{-20pt}
\begin{itemize}\setlength\itemsep{3pt}
  \item Two variables likely endogenous: Adjust for changes induced by dismissals
  \begin{itemize}\setlength\itemsep{3pt}
    \item Dismissal induced reduction in faculty quality (higher value $\to$ larger fall)\vspace{-5pt}
    \[ (Avg. \text{1932 Quality})-(Avg. \text{1932 Quality of stayers})
    \]
    \item Dismissal induced increase in S/F ratio: 
    \[\frac{\# \text{1932 student}}{\# \text{1932 Prof}-\# \text{Dismissed Prof}}-\frac{\# \text{1932 student}}{\# \text{1932 Prof}}\]
  \end{itemize}
\end{itemize}
\end{frame}

\begin{frame}
\frametitle{Significant drop-off in quality of departments with dismissals} 
\begin{figure}
  \begin{subfigure}{0.49\textwidth}
    \includegraphics[keepaspectratio,width=\textwidth]{w2010_f1.png}
    \caption{Department size}
  \end{subfigure}
  \begin{subfigure}{0.49\textwidth}
    \includegraphics[keepaspectratio,width=\textwidth]{w2010_f1_2.png}
    \caption{Faculty quality}
  \end{subfigure}
\end{figure}
\end{frame}


\begin{frame}
\frametitle{Drop-off in PhD student outcomes correlated with quality} 
\begin{figure}
    \includegraphics[keepaspectratio,width=0.9\textwidth]{w2010_t1.png}
\end{figure}
\end{frame}

\begin{frame}
\frametitle{Effects driven by changes in top 10 institutions} 
\begin{figure}
  \begin{subfigure}{0.49\textwidth}
    \includegraphics[keepaspectratio,width=\textwidth]{w2010_t2_1.png}
  \end{subfigure}
  \begin{subfigure}{0.49\textwidth}
    \includegraphics[keepaspectratio,width=\textwidth]{w2010_t2_2.png}

  \end{subfigure}
\end{figure}
\end{frame}

\begin{frame}
\frametitle{Takeaways and discussions}
Departure of talent does affect productivity and outcomes of remaining students 
\begin{itemize}\setlength\itemsep{3pt}
\item Quality of faculty has sizeable effects on student outcomes at German Universities
\item Quality of instruction matters even in selective settings
\item Departure of key talents could have long-run productivity consequence at home
\end{itemize}\medskip
Remaining questions
\begin{itemize}\setlength\itemsep{3pt}
\item Are these effects persistent? Depends on whether the talent returns home
\begin{itemize}\setlength\itemsep{3pt}
\item Not applicable here: These professors could not return to Germany
\item Return migrants can transfer knowledge and resources from foreign countries 
\end{itemize}
\item Other ways in which migrants affect innovation and productivity
\begin{itemize}\setlength\itemsep{3pt}
\item Studies using how visa restrictions (H1B in US)$\to$ $\Downarrow$ patents, $\Uparrow$ Outsourcing by US firms
\item Potentially suggest that migrant labor may not be replaced with domestic labor
\end{itemize}
\end{itemize}
\end{frame}

\subsection{Mahajan, Yang (2020): Taken by Storm: Hurricanes, Migrant Networks, and US Immigration}
\begin{transitionframe}
\begin{center}
  \begin{huge}
    \textcolor{NTHU1}{%
  Mahajan, Yang (2020):\\ Taken by Storm: Hurricanes, Migrant Networks, and US Immigration
    }
  \end{huge}
\end{center}
\end{transitionframe}

\begin{frame}
\frametitle{Q: Do negative shocks at origin countries increase migration?} 
Many unanswered questions in migration, with mixed theoretical and empirical findings
\begin{itemize}\setlength\itemsep{3pt}
\item Negative shocks at home lead to increase or decrease in international migration?
\item What roles do network of prior migrants play in outmigration decision?
\item Are these driven by legal or illegal channels?
\item Are these migrations likely to be temporary or permanent?
\end{itemize}\medskip
This paper uses hurricanes to estimate impact of natural disasters on migration decisions
\begin{itemize}\setlength\itemsep{3pt}
\item 25 years of hurricanes data covering 159 origin countries of US migrants
\item Examines how prior migrant presence affects future migrant flows
\item Uses these data to answer theoretical ambiguities
\begin{itemize}\setlength\itemsep{3pt}
  \item Individuals decide to stay home vs move to attractive location
  \item Rising fixed cost of moving out vs Returns to migration higher
  \item Migrant networks may reduce migration costs vs reduce migration needs ($\because$ remittances)
\end{itemize}
\end{itemize}\medskip
\end{frame}

\begin{frame}
\frametitle{Conceptual justification can go either way} 
Stay at home location vs move to more attractive location
\begin{itemize}\setlength\itemsep{3pt}
\item If one chooses to move, there is a fixed cost to migrate
\item After shock: Destination becomes (relatively) more attractive vs Moving costs rises
\begin{itemize}\setlength\itemsep{3pt}
\item Increased demand raises cost of migration services
\item Make it difficult to obtain credit to finance migration
\end{itemize}
\end{itemize}\medskip
Migrant network may also have a complicated role
\begin{itemize}\setlength\itemsep{3pt}
\item Earlier research: Larger network$\to$ $\Downarrow$ Fixed costs
\begin{itemize}
\item Reduce search and information costs, social support, and sponsor relatives 
\end{itemize}
\item Prior migrants as insurance network, reducing need to migrate
\begin{itemize}\setlength\itemsep{3pt}
\item Can send financial assistance home (remittances) that helps with coping the shock
\item This can happen instead of helping them migrate in response to shocks
\end{itemize}
\end{itemize}\medskip
$\Rightarrow$Uses empirical analysis to identify which channel is working
\end{frame}


\begin{frame}
\frametitle{Data and empirical strategy} 
Data sources
\begin{itemize}\setlength\itemsep{3pt}
\item US immigration rates: 1980-2004 for each observable origins (US Census data)
\item Hurricanes: Tracks from National Oceanic and Atmospheric Administration
\item Migrant networks: Census, American Community Survey, and DHS data
\end{itemize}\medskip
Leverage variation in returns to migration across time and origin locations
\[
Y_{jt} = \beta_0 + \beta_1H_{jt}+\beta_2(H_{jt}\times s_{j,1980})+\eta_j+\delta_t + \phi_j\times t + e_{jt}
\]\vspace{-20pt}
\begin{itemize}\setlength\itemsep{3pt}
\item $H_{jt}$: Hurricane index for origin $j$ at year $t$
\item $s_{j,1980}$: Stock of migrants from $j$ at US in 1980 
\item $\beta_1$ measures effect of hurricanes on migration 
\item $\beta_2$ captures differentiated effects of hurricanes based on prior migrant network 
\end{itemize}\medskip
\end{frame}

\begin{frame}
\frametitle{How does the hurricane track data look like?} 
\begin{figure}
    \includegraphics[keepaspectratio,width=0.8\textwidth]{my2020_f1.png}
\end{figure}
\end{frame}

\begin{frame}
\frametitle{Increased migration, amplified by presence of migrant networks} 
\begin{figure}
    \includegraphics[keepaspectratio,width=0.8\textwidth]{my2020_t1.png}
\end{figure}
\end{frame}

\begin{frame}
\frametitle{Increases are driven by naturalized migrants} 
\begin{figure}
    \includegraphics[keepaspectratio,width=0.8\textwidth]{my2020_t2.png}
\end{figure}
\end{frame}

\begin{frame}
\frametitle{Additional findings and remaining questions}
Additional findings 
\begin{itemize}\setlength\itemsep{3pt}
\item Using different sources of migration data (DHS) does not change results
\item Effects are driven by youngest (0-12 y.o) and prime age (18-44 y.o) migrants
\begin{itemize}\setlength\itemsep{3pt}
\item Those who migrate for work related reasons and their children
\end{itemize}
\item Family sponsorships (not refugees or employment sponsor) are key networks
\item This is driven by increase in desire to migrate
\begin{itemize}\setlength\itemsep{3pt}
\item Leniency programs by US (TPS and DED) did not change
\end{itemize}
\end{itemize}\medskip
Remaining questions
\begin{itemize}\setlength\itemsep{3pt}
\item Implications for climate-driven migration
\begin{itemize}
\item Increasing importance due to increases in natural disasters
\end{itemize}
\item Causation for any studies with immigration treatment: Selection due to prior networks
\begin{itemize}
\item Shift-share IV: Early-stage migrants shares and contemporary shifts
\end{itemize}
\end{itemize}\medskip
\end{frame}

\subsection{Becker, Grosfeld, Grosjean, Voigländer, and Zhuravskaya (2020): Forced Migration and Human Capital: Evidence from Post WWII Population Transfers [Student presentation]}
\begin{transitionframe}
\begin{center}
  \begin{huge}
    \textcolor{NTHU1}{%
Becker, Grosfeld, Grosjean, Voigländer, and Zhuravskaya (2020):\\ Forced Migration and Human Capital: Evidence from Post WWII Population Transfers [Student presentation]
    }
  \end{huge}
\end{center}
\end{transitionframe}





\section{How do natives (hosts) perceive migrants}
\subsection{Fouka, Tabellini (2022): Changing In-Group Boundaries: The Effect of Immigration on Race Relations in the United States}
\begin{transitionframe}
\begin{center}
  \begin{huge}
    \textcolor{NTHU1}{%
Fouka, Tabellini (2022): \\ Changing In-Group Boundaries: The Effect of Immigration on Race Relations in the United States
    }
  \end{huge}
\end{center}
\end{transitionframe}

\begin{frame}
\frametitle{Q: How are in/out-group boundaries formed in multigroup relations?} 
Categorization of in/out-group members explain attitudes of the majority group
\begin{itemize}\setlength\itemsep{3pt}
\item In-group favoritism and out-group prejudice $\to$ Possible conflict \& violence
\item The boundaries change through interactions $\to$ Even in multigroup setting?
\item Demographic shifts from migration serves as natural experiment
\begin{itemize}\setlength\itemsep{3pt}
\item Introduction of a new minority group
\item Shift prejudice from one minority to other? Lump all minorities together?
\end{itemize}
\end{itemize}\medskip
This paper studies how inflow of new migrants affects perception of existing minorities
\begin{itemize}\setlength\itemsep{3pt}
\item Leverage variation in Mexican migrants and perception on Black Americans
\item \textit{Affective distance} and size determines which attributes matter for social categorization
\begin{itemize}\setlength\itemsep{3pt}
\item Individual's feelings towards members of different groups relative to their own
\item Relative size affects intergroup relations
\end{itemize}
\end{itemize}
\end{frame}


\begin{frame}
\frametitle{Conceptual Framework: Dynamics of categorizing  in and out-groups}

Attributes used for classification are determined by group share and affective distance
\begin{itemize}\setlength\itemsep{3pt}
\item Classification maximizes ratio of across-group differences over in-group differences
\item For attribute $j$, let $I_j^k=1$ if group $k$ differ from in-group in terms of $j$ (0 if same)
\item For group $k$ with size $n^k$ with affective distance $\delta^k$, salient attribute $j$ maximizes
\[
\max_j R_j =\left. \underbrace{\frac{\sum_k \delta^k n^k I_j^k}{\sum_k n^k I_j^k}}_{\substack{\text{weighted avg of} \\ \delta\text{ across out-groups}
}}\middle/ \underbrace{\frac{\sum_k \delta^k n^k (1-I_j^k)}{\sum_k n^k (1-I_j^k)}}_{\substack{\text{weighted avg of} \\ \delta\text{ across in-groups}
}}\right.
\]
\end{itemize}\medskip \vspace{-10pt}
Takeaway % Refer to page 4 (proposition 2)
\begin{itemize}\setlength\itemsep{3pt}
\item Group composition change (size) $\to$ Salient attribute changes (immigration over race)
\item Existing groups can be recategorized from outsider to insider depending on
\begin{itemize}\setlength\itemsep{3pt}
\item Whether it differs from rising groups in terms of attributes
\item Affective distance and size of the incoming group
\end{itemize}
\end{itemize}
\end{frame}

\begin{frame}
\frametitle{Data and empirical strategy}

Data sources
\begin{itemize}\setlength\itemsep{3pt}
\item Migrants: Census (1970-2010), American National Election Study (ANES)
\item Sentiment: Feelings thermometer + stereotypical views from ANES 

\end{itemize}\medskip 
Empirical strategy leverages Mexican and other migrant inflow and attitudes (DID)
\[
y_{irst}=\beta_1 M_{rst} + \beta_2 S_{rst} + X_{irst} + \gamma_{rs}+\mu_{rt} + \eta_{irst}
\]\vspace{-20pt}
\begin{itemize}
\item $M_{rst}$: Share of Mexican migrants in census tract $r$ in state $s$ at time $t$
\item $S_{rst}$: Share of non-Mexican migrants in census tract $r$ in state $s$ at time $t$
\item $M_{rst}$ and $S_{rst}$ can be endogenous \medskip
\begin{itemize}\setlength\itemsep{3pt}
\item Instrument with predicted share of migrants using 1960 settlement shares (Shift-share IV)
\item Relevance: Idea is that new inflow of migrants tends to follow previous ones 
\item Exogeneity: States with high share of Mexicans in `60 not on different $y_{irst}$ trajectory
\end{itemize}
\end{itemize}
\end{frame}





\begin{frame}
\frametitle{Favorable towards Blacks, not so for Hispanics}
\begin{center}
\begin{figure}
\begin{subfigure}{0.75\textwidth}
\includegraphics[width=\textwidth]{ft2022_t1.jpg}
\end{subfigure}
\begin{subfigure}{0.75\textwidth}
\includegraphics[width=\textwidth]{ft2022_t2.jpg}
\end{subfigure}
\end{figure}
\end{center}

\end{frame}

\begin{frame}
\frametitle{Driven by increased salience of immigration over race}
\begin{center}
\begin{figure}
\includegraphics[width=0.9\textwidth, keepaspectratio]{ft2022_t3.jpeg}
\end{figure}
\end{center}
\end{frame}

\begin{frame}
\frametitle{Composition of hate crime victims change}
\begin{center}
\begin{figure}
\includegraphics[width=0.725\textwidth, keepaspectratio]{ft2022_t4.png}
\end{figure}
\end{center}
\end{frame}



\begin{frame}
\frametitle{Takeaways and discussions} % Why I need to test this

Inflow of migration changes boundaries of who is in or out 
\begin{itemize}\setlength\itemsep{3pt}
\item Distant group size (Mexicans) $\Uparrow$ $\to$ favorable attitudes towards existing minorities
\item Shifting salient issue (migration) determines recategorization of in/out-groups
\item Results generalized to wider set of ethnicities (Asians in-group, Muslims out-group)
\end{itemize}\medskip
Remaining questions
\begin{itemize}\setlength\itemsep{3pt}
\item Is size the most pressing factor in determining affective distance?
\begin{itemize}\setlength\itemsep{3pt}
\item Southeast Asian worker migrants still viewed out-group, despite size
\item Other factors contribute to affective distance? (e.g. economic/educational background?)
\end{itemize}
\item Are all ingroup members the same?
\begin{itemize}\setlength\itemsep{3pt}
\item Black Americans still not considered equal to White Americans
\item Differences within Hispanics (Cuban and Venezuelans vs other Hispanics)
\item In-groups can be broken up? (e.g. media / socioeconomics / gentrification)
\end{itemize}
\end{itemize}
\end{frame}

\subsection{Alesina, Miano, Stantcheva (2023): Immigration and Redistribution}
\begin{transitionframe}
\begin{center}
  \begin{huge}
    \textcolor{NTHU1}{%
  Alesina, Miano, Stantcheva (2023):\\ Immigration and Redistribution
    }
  \end{huge}
\end{center}
\end{transitionframe}

\begin{frame}
\frametitle{Q: Does immigration change policy support?}
Conflicts on policy designs, particularly redistribution,  following increase in immigration
\begin{itemize}\setlength\itemsep{3pt}
\item Do citizens have accurate perception about migrants?
\begin{itemize}\setlength\itemsep{3pt}
  \item Do they over/underestimate the share of migrant population?
  \item Do they believe that immigrants burden welfare system without contributing to it?
  \item Do they believe that immigrants are culturally distant?
\end{itemize}
\item What is the link between the perceptions and support for redistribution policy?
\begin{itemize}\setlength\itemsep{3pt}
  \item Do treatments that aim to correct perceptions affect opinions on these policies?
\end{itemize}
\end{itemize}\medskip
This paper uses survey experiment to address how immigration affects policy preferences
\begin{itemize}\setlength\itemsep{3pt}
  \item Online survey on 22,000 people across 6 countries (US, UK, FRA, GER, ITA, SWE)
  \item Gather information on perception on migrants and preferences for redistribution
  \item Various treatment to test if highlighting migration affects policy opinions
\end{itemize}
\end{frame}

\begin{frame}
\frametitle{Data and empirical strategy}
Randomized survey experiment design: salience treatment and information treatment
\begin{itemize}\setlength\itemsep{3pt}
  \item Background info: Gender, age, education, employment, political orientation etc
\item Salience treatment: Immigration block before redistribution block (control: vice versa)
\begin{itemize}\setlength\itemsep{3pt}
\item \textit{Tests whether making migration salient affects policy preferences}
\item Immigration: Ask perception on immigrants, views on immigration policy
\item Redistribution: Ask redistribution and govt role questions (no reference to migrants)
\end{itemize}
\item Information treatment: Control group and three types of information
\begin{itemize}\setlength\itemsep{3pt}
\item \textit{Tests whether providing additional information changes the effect of salience }
\item Video 1 on the share of immigrants
\item Video 2 on the origins of the immigrants
\item Video 3 that emphasizes hard work of immigrants (no facts)
\end{itemize}
\item 8 groups of treatment arms: Uses OLS to compare across the groups
\end{itemize}\medskip
\end{frame}

\begin{frame}
\frametitle{Content of the video experiment}
\begin{figure}
  \includegraphics[keepaspectratio, width=0.6\textwidth]{ams2023_f1.png}
\end{figure}
\end{frame}

\begin{frame}
\frametitle{There are drastic misperception on migrants}
\begin{figure}
  \begin{subfigure}{0.35\textwidth}
    \includegraphics[keepaspectratio, width=\textwidth]{ams2023_f2_1.png}
    \caption{Share of immigrants}
  \end{subfigure}
  \begin{subfigure}{0.35\textwidth}
    \includegraphics[keepaspectratio, width=\textwidth]{ams2923_f2_2.png}
    \caption{Share of Muslim immigrants}
  \end{subfigure} 
\end{figure}
Note: Also documented misperception on education, attribute of poverty status
\end{frame}

\begin{frame}
\frametitle{Salience reduces support, limited correction by information}
\begin{figure}
  \includegraphics[keepaspectratio, width=0.9\textwidth]{ams2023_t1.png}
\end{figure}
\end{frame}

\begin{frame}
\frametitle{Driven by misperception on economic status of immigrants}
\begin{figure}
  \includegraphics[keepaspectratio, width=0.9\textwidth]{ams2023_f3.png}
\end{figure}
\end{frame}

\begin{frame}
\frametitle{Takeaways and discussions}
Much of the debates are actually based on misperceptions on immigrants
\begin{itemize}\setlength\itemsep{3pt}
\item Overestimate size, cultural distance, and poverty status of migrants
\item Just making people think about migration skews policy preferences
\begin{itemize}\setlength\itemsep{3pt}
\item Driven by misperception on economic status of immigrants
\item Particularly dominant on older, less educated, poorer, and right-wing respondents
\end{itemize}
\item Follow-up: Those with larger misperceptions less willing to pay for correct information
\end{itemize}\medskip
Remaining questions on policy preferences
\begin{itemize}\setlength\itemsep{3pt}
\item How to correct misperceptions that can skew policy preferences?
\begin{itemize}\setlength\itemsep{3pt}
\item Otherwise, could be used by anti-immigration parties (or extremists)
\item Worrisome given that those who need this more are less willing
\end{itemize}
\end{itemize}
\end{frame}

\subsection{Hainmueller, Hiscox (2010): Attitudes toward Highly-Skilled and Low-Skilled Immigration: Evidence from a Survey Experiment [student presentation]}
\begin{transitionframe}
\begin{center}
  \begin{huge}
    \textcolor{NTHU1}{%
Hainmueller, Hiscox (2010):\\ Attitudes toward Highly-Skilled and Low-Skilled Immigration: Evidence from a Survey Experiment [student presentation]
    }
  \end{huge}
\end{center}
\end{transitionframe}

\begin{frame}
\frametitle{Takeaways and remaining questions} 
Migrants do contribute, but are often misunderstood
\begin{itemize}\setlength\itemsep{3pt}
\item Contribute in labor market, human capital, and politics
\item Yet misperceptions about their efforts inhibit integration and further contribution
\begin{itemize}\setlength\itemsep{3pt}
 \item Natives often consider them as out-group and view them with prejudice
\item Often perpetuated by lack of credible information and misinformation from media
\end{itemize}
\item Often leads to further social and political conflicts
\begin{itemize}\setlength\itemsep{3pt}
 \item Policy decisions are skewed against their favor
\item Often victimized by crime and conflict even in host societies
\end{itemize}
\end{itemize}\medskip
Remaining questions
\begin{itemize}\setlength\itemsep{3pt}
 \item How do you improve two-way efforts to increase migrant-native integration?
 \begin{itemize}\setlength\itemsep{3pt}
 \item Ways to address economic/cultural barriers and misinformation?
\end{itemize}
\item Studies on political refugees and climate migrants
\begin{itemize}\setlength\itemsep{3pt}
 \item Limited possibility to return to homeland $\to$ Understanding integration is crucial
\end{itemize}
\end{itemize}
\end{frame}


%%%%%%%%%%%
\end{document}


